\begin{frame}{讲解：实验方法课的中心不是跑图，而是写出可否定命题}
  \Core{如果实验只是为了证明自己是对的，那课程项目和论文都很容易滑向结果导向叙事。}
  \begin{itemize}
    \item 好实验首先要写清假设，再写清什么结果会直接否定这个假设。
    \item workload、对照组和探针的设计，本质上都是围绕这条可否定命题服务的。
    \item 因此实验方法课训练的是研究纪律，而不是仅仅训练脚本使用能力。
  \end{itemize}
\end{frame}

\begin{frame}{讲解：端到端指标与局部探针为什么必须成对出现}
  \Core{只报端到端结果，结论容易悬空；只报局部探针，结果又容易失去服务意义。}
  \begin{itemize}
    \item 端到端指标告诉我们用户看到了什么，例如 TTFT、P99、goodput。
    \item 局部探针告诉我们系统里发生了什么，例如 batch 变化、KV 占用、编译命中。
    \item 两者结合起来，才能把“现象”变成“因果解释”。
  \end{itemize}
\end{frame}

\begin{frame}{例子：一个看起来完整但其实不可信的实验（1）}
  \Core{实验方法课必须先教学生识别坏实验，否则工具越熟练，错误可能越系统化。}
\begin{block}{表面上很完整}
    有吞吐图、有延迟表、甚至有消融，但 workload 来源和对照配置没有交代清楚。
  \end{block}
\end{frame}

\begin{frame}{例子：一个看起来完整但其实不可信的实验（2）}
  \Core{实验方法课必须先教学生识别坏实验，否则工具越熟练，错误可能越系统化。}
\begin{block}{真正的问题}
    结果无法回答“为什么变好”以及“在哪些条件下不再成立”。
  \end{block}
\end{frame}

\begin{frame}{例子：一个看起来完整但其实不可信的实验（3）}
  \Core{实验方法课必须先教学生识别坏实验，否则工具越熟练，错误可能越系统化。}
\begin{block}{课程结论}
    一个实验是否可信，不取决于图表数量，而取决于因果链和边界是否被写清。
  \end{block}
\end{frame}

\begin{frame}{例子：端到端数字与探针数字不一致时怎么办（1）}
  \Core{最有研究价值的情况往往不是数字全都一致，而是局部证据与端到端现象出现张力。}
\begin{block}{情况一}
    GPU 利用率更高，但 TTFT 更差，说明等待或同步代价可能在放大。
  \end{block}
\end{frame}

\begin{frame}{例子：端到端数字与探针数字不一致时怎么办（2）}
  \Core{最有研究价值的情况往往不是数字全都一致，而是局部证据与端到端现象出现张力。}
\begin{block}{情况二}
    KV 命中率更高，但 goodput 没提升，说明命中本身未必落在主导瓶颈上。
  \end{block}
\end{frame}

\begin{frame}{例子：端到端数字与探针数字不一致时怎么办（3）}
  \Core{最有研究价值的情况往往不是数字全都一致，而是局部证据与端到端现象出现张力。}
\begin{block}{教学意义}
    探针不是为了给端到端结果背书，而是为了帮助发现矛盾并修正假设。
  \end{block}
\end{frame}

\begin{frame}{代码例子：最小 benchmark 模板应包含什么（1）}
  \Core{一份合格的 benchmark，不只是能跑，还应显式携带 workload 和环境信息。}
\begin{block}{代码片段}
    {\ttfamily\small
    \detokenize{seed(0)}\par
    \detokenize{warmup_requests = 10}\par
    \detokenize{num_requests = 100}\par
    \detokenize{measure(ttft, throughput, p99)}\par
    }
  \end{block}
\end{frame}

\begin{frame}{代码例子：最小 benchmark 模板应包含什么（2）}
  \Core{一份合格的 benchmark，不只是能跑，还应显式携带 workload 和环境信息。}
\begin{block}{课堂应追问}
    哪些配置是影响 workload 定义的，哪些配置只是运行时细节？
  \end{block}
\end{frame}

\begin{frame}{代码例子：最小 benchmark 模板应包含什么（3）}
  \Core{一份合格的 benchmark，不只是能跑，还应显式携带 workload 和环境信息。}
\begin{block}{系统含义}
    benchmark 模板的显式化，本身就在减少“实验只存在于一台机器上”的偶然性。
  \end{block}
\end{frame}

\begin{frame}{代码例子：最小日志汇总脚本（1）}
  \Core{实验课需要让学生看到，图表往往来自对原始日志的二次组织。}
\begin{block}{代码片段}
    {\ttfamily\small
    \detokenize{records = load_jsonl("run.log") }\par
    \detokenize{p99 = percentile([r["latency_ms"] for r in records], 99)}\par
    \detokenize{goodput = sum(r["ok"] for r in records) / total_time}\par
    }
  \end{block}
\end{frame}

\begin{frame}{代码例子：最小日志汇总脚本（2）}
  \Core{实验课需要让学生看到，图表往往来自对原始日志的二次组织。}
\begin{block}{课堂应追问}
    如果日志字段设计不对，后面哪些指标将根本算不出来？
  \end{block}
\end{frame}

\begin{frame}{代码例子：最小日志汇总脚本（3）}
  \Core{实验课需要让学生看到，图表往往来自对原始日志的二次组织。}
\begin{block}{课程结论}
    可验证性不只发生在论文写作阶段，而是从日志字段设计就已经开始了。
  \end{block}
\end{frame}

\begin{frame}{例子：消融实验最容易做假的地方（1）}
  \Core{消融的真正任务是隔离变量，而不是制造更多子图。}
\begin{block}{伪消融一}
    一次改了多个控制点，却只关掉其中一个开关就说完成了消融。
  \end{block}
\end{frame}

\begin{frame}{例子：消融实验最容易做假的地方（2）}
  \Core{消融的真正任务是隔离变量，而不是制造更多子图。}
\begin{block}{伪消融二}
    workload 或参数同时变化，导致数字变化无法归因。
  \end{block}
\end{frame}

\begin{frame}{例子：消融实验最容易做假的地方（3）}
  \Core{消融的真正任务是隔离变量，而不是制造更多子图。}
\begin{block}{课程要求}
    每一项消融都要明确写出“被隔离的变量”和“预期影响的指标”。
  \end{block}
\end{frame}

\begin{frame}{研究案例：shared workload 不是“拿来就能跑”}
  \Core{真实 benchmark 的 provenance 包括 workload 来源，也包括 adapter 如何把它变成当前系统可运行的输入。}
  \begin{itemize}
    \item KV-Delta 中，shared workload 元数据里的 prompt length 并不等于 Qwen tokenizer token 数。
    \item 原始请求真实 token 数超过 \texttt{\detokenize{max_model_len}}，导致 completion endpoint 返回 400。
    \item 正确做法是记录 token-budget 裁剪、deterministic subset、submodule commit 和失败 artifact。
  \end{itemize}
\end{frame}

\begin{frame}{研究案例：同一实验报告必须保留失败 run}
  \Core{failed / blocked run 不是噪音，而是解释系统边界和后续工作的证据。}
  \begin{itemize}
    \item full 40-request shared case 能通过请求阶段，但 consumer trace/save-load 后段过重。
    \item 因此改用 deterministic first-8 subset 完成 m=2/4/8 canonical，不等于完整 workload 已经解决。
    \item 课程要求：把 full-40 标为 blocked，把 first-8 标为 shared subset，不得混写。
  \end{itemize}
\end{frame}

\begin{frame}{研究案例：claim ledger 的最小字段}
  \Core{每个课程项目都应把结果写成可审计 claim，而不是只贴一张性能图。}
  \begin{itemize}
    \item Claim：当前结论能说到什么程度。
    \item Evidence：artifact path、workload source、baseline、指标和设备。
    \item Boundary：invalid/failed/blocked run、未覆盖 workload、尚未解释的反例。
  \end{itemize}
\end{frame}

\begin{frame}{课堂练习：写一页可信实验卡片（1）}
  \Core{实验方法训练最好用一页纸完成：问题、对照、指标、失败条件。}
\begin{block}{练习一}
    写出你的 workload 定义，至少包含输入长度分布、输出长度分布和到达模式。
  \end{block}
\end{frame}

\begin{frame}{课堂练习：写一页可信实验卡片（2）}
  \Core{实验方法训练最好用一页纸完成：问题、对照、指标、失败条件。}
\begin{block}{练习二}
    写出一个主指标和两个辅助探针，并说明三者的因果关系假设。
  \end{block}
\end{frame}

\begin{frame}{课堂练习：写一页可信实验卡片（3）}
  \Core{实验方法训练最好用一页纸完成：问题、对照、指标、失败条件。}
\begin{block}{练习三}
    用一句话写出：什么结果出现时，你会承认当前机制没有兑现预期收益？
  \end{block}
\end{frame}
